\documentclass[11pt,a4paper]{article}

% ── Packages ──────────────────────────────────────────────────────────────────
\usepackage[margin=2.5cm]{geometry}
\usepackage[T1]{fontenc}
\usepackage{lmodern}
\usepackage{xcolor}
\usepackage{booktabs}
\usepackage{longtable}
\usepackage{tabularx}
\usepackage{multirow}
\usepackage{enumitem}
\usepackage{fancyhdr}
\usepackage{titlesec}
\usepackage{hyperref}
\usepackage{listings}
\usepackage{mdframed}
\usepackage{graphicx}
\usepackage{array}
\usepackage{colortbl}
\usepackage{parskip}

% ── Colours ───────────────────────────────────────────────────────────────────
\definecolor{netframe}{RGB}{20,30,60}
\definecolor{critical}{RGB}{180,20,20}
\definecolor{high}{RGB}{200,80,0}
\definecolor{medium}{RGB}{180,140,0}
\definecolor{low}{RGB}{40,130,40}
\definecolor{resolved}{RGB}{0,120,60}
\definecolor{accepted}{RGB}{80,80,160}
\definecolor{partial}{RGB}{160,100,0}
\definecolor{pending}{RGB}{120,60,0}
\definecolor{codebg}{RGB}{245,245,245}
\definecolor{rowalt}{RGB}{248,250,255}

% ── Page style ────────────────────────────────────────────────────────────────
\pagestyle{fancy}
\fancyhf{}
\renewcommand{\headrulewidth}{0.4pt}
\fancyhead[L]{\small\color{netframe}\textbf{NetFRAME Homelab}}
\fancyhead[R]{\small\color{netframe}Penetration Test Remediation Report}
\fancyfoot[C]{\small\color{gray}\thepage}
\fancyfoot[R]{\small\color{gray}CONFIDENTIAL}

% ── Section formatting ────────────────────────────────────────────────────────
\titleformat{\section}{\large\bfseries\color{netframe}}{}{0em}{}[\titlerule]
\titleformat{\subsection}{\normalsize\bfseries\color{netframe!80}}{}{0em}{}
\titleformat{\subsubsection}{\normalsize\bfseries}{}{0em}{}

% ── Code listings ─────────────────────────────────────────────────────────────
\lstset{
  basicstyle=\small\ttfamily,
  backgroundcolor=\color{codebg},
  frame=single,
  framerule=0.4pt,
  rulecolor=\color{gray!40},
  breaklines=true,
  breakatwhitespace=true,
  xleftmargin=6pt,
  xrightmargin=6pt,
  aboveskip=6pt,
  belowskip=6pt,
}

% ── Severity badge macros ─────────────────────────────────────────────────────
\newcommand{\critical}{\colorbox{critical}{\color{white}\small\textbf{CRITICAL}}}
\newcommand{\high}{\colorbox{high}{\color{white}\small\textbf{HIGH}}}
\newcommand{\medium}{\colorbox{medium}{\color{white}\small\textbf{MEDIUM}}}
\newcommand{\low}{\colorbox{low}{\color{white}\small\textbf{LOW}}}
\newcommand{\resolved}{\colorbox{resolved}{\color{white}\small\textbf{RESOLVED}}}
\newcommand{\riskaccepted}{\colorbox{accepted}{\color{white}\small\textbf{RISK ACCEPTED}}}
\newcommand{\partialstatus}{\colorbox{partial}{\color{white}\small\textbf{PARTIAL}}}
\newcommand{\pendingvlan}{\colorbox{pending}{\color{white}\small\textbf{PENDING VLAN}}}

% ─────────────────────────────────────────────────────────────────────────────
\begin{document}

% ── Cover page ────────────────────────────────────────────────────────────────
\begin{titlepage}
  \centering
  \vspace*{3cm}

  {\Huge\bfseries\color{netframe} NetFRAME Homelab}\\[0.4em]
  {\LARGE\color{netframe!70} Penetration Test Remediation Report}\\[2em]

  \rule{\linewidth}{1.5pt}\\[1em]

  \begin{tabular}{ll}
    \textbf{Assessment Date:} & 2026-06-23 \\[0.3em]
    \textbf{Remediation Date:} & 2026-06-24 \\[0.3em]
    \textbf{Scope:} & 192.168.10.0/24 (Management VLAN) \\[0.3em]
    \textbf{Client:} & Kyle Mason \\[0.3em]
    \textbf{Environment:} & NetFRAME CS9000 42U Rack, Greater Cleveland OH \\[0.3em]
    \textbf{Classification:} & CONFIDENTIAL \\[0.3em]
    \textbf{Report Version:} & 1.0 \\
  \end{tabular}

  \vfill

  \rule{\linewidth}{0.4pt}\\[0.5em]
  {\small This report documents the findings and remediation actions taken following an internal
  penetration test of the NetFRAME management network. All testing and remediation was
  conducted on systems owned and operated by the report recipient.}

\end{titlepage}

% ── Table of Contents ─────────────────────────────────────────────────────────
\tableofcontents
\newpage

% ─────────────────────────────────────────────────────────────────────────────
\section{Executive Summary}

An internal penetration test of the NetFRAME homelab management network
(\texttt{192.168.10.0/24}) was conducted on 2026-06-23, followed by a structured remediation
engagement on 2026-06-24. The assessment identified \textbf{10 findings} across the
management VLAN, ranging from critical power-infrastructure exposure to IoT devices on the
wrong network segment.

\textbf{Key results:}
\begin{itemize}[nosep]
  \item 34 live hosts discovered; 11 previously undocumented
  \item 5 findings \textbf{fully resolved} during this engagement
  \item 1 finding \textbf{risk accepted} (low severity, managed device)
  \item 1 finding \textbf{partially remediated} (pending hardware arrival)
  \item 3 findings \textbf{identified and documented} (pending VLAN segmentation)
  \item No critical data accessed; no destructive actions taken
\end{itemize}

\textbf{Most significant finding:} The Middle Atlantic UPS-OL2200R management card (which
controls power to the compute stack) exposed Telnet and FTP — cleartext protocols over
power-critical infrastructure. This was fully remediated.

\textbf{Positive observations:} OPNsense firewall operational, Proxmox Backup Server live with
daily backup schedules, Headscale VPN deployed, Vaultwarden operational, Grafana requires
authentication, NPM default credentials not in use.

% ─────────────────────────────────────────────────────────────────────────────
\section{Findings Summary}

\begin{longtable}{@{}p{1cm}p{5.5cm}p{2.5cm}p{2cm}p{2cm}@{}}
\toprule
\textbf{ID} & \textbf{Title} & \textbf{Severity} & \textbf{Status} & \textbf{Date} \\
\midrule
\endhead
F-01 & Undocumented Standalone Node (pve1) & \medium & \partialstatus & 2026-06-23/24 \\[3pt]
F-02 & UPS Management Card — Telnet/FTP Exposed & \critical & \resolved & 2026-06-24 \\[3pt]
F-03 & Prometheus API Unauthenticated & \high & \resolved & 2026-06-24 \\[3pt]
F-04 & Homepage Dashboard — No Authentication & \high & \resolved & 2026-06-24 \\[3pt]
F-05 & NPM Admin Panel Exposed on LAN & \medium & \resolved & 2026-06-24 \\[3pt]
F-06 & Weak 1024-bit RSA SSH Host Key (.176) & \medium & \riskaccepted & 2026-06-24 \\[3pt]
F-07 & Amazon Fire Tablet on Management VLAN & \medium & \pendingvlan & 2026-06-24 \\[3pt]
F-08 & Duplicate Pi-hole — DNS Conflict Risk & \medium & \resolved & 2026-06-23 \\[3pt]
F-09 & Unidentified Devices .104---.108 & \low & \pendingvlan & 2026-06-24 \\[3pt]
F-10 & Unidentified Device BBL-P003 (.110) & \low & \pendingvlan & 2026-06-24 \\
\bottomrule
\end{longtable}

\vspace{0.5em}
\noindent\textbf{Note on F-08:} Closed as a false positive during the assessment. The
Pi-hole at \texttt{192.168.10.177} is hosted exclusively on pve1 (Mac Mini LXC 103); pve3
has no Pi-hole instance. The original finding was based on incorrect documentation.

% ─────────────────────────────────────────────────────────────────────────────
\section{Detailed Findings and Remediation}

% ── F-01 ─────────────────────────────────────────────────────────────────────
\subsection{F-01 — Undocumented Standalone Node (pve1) \quad \medium \quad \partialstatus}

\textbf{Host:} \texttt{192.168.10.193} (Mac Mini)

\textbf{Description:} A Mac Mini running Proxmox VE 9.1.9 was found active on the management
network, not present in any asset inventory. The node was intentionally removed from the
km-cluster (Broadcom NIC incompatible with corosync knet) but retained to serve two active
LXC containers: Pi-hole DNS (LXC 103, \texttt{192.168.10.177}) and a Homepage dashboard
(LXC 104, \texttt{192.168.10.174}). The Pi-hole at \texttt{.177} is the primary DNS server
for the entire network — shutdown without migration would cause full DNS failure.

\textbf{Risk:} Not enrolled in Wazuh SIEM; not covered by PBS backups; shared SSH
\texttt{authorized\_keys} with cluster nodes; Homepage dashboard accessible without
authentication.

\textbf{Remediation Applied:}
\begin{itemize}[nosep]
  \item \textbf{SSH hardened:} \texttt{PermitRootLogin prohibit-password} set; 5 stale cluster RSA
    keys removed from \texttt{/root/.ssh/authorized\_keys}; only Ares Ed25519 key retained.
  \item \textbf{Wazuh enrolled:} Agent 4.9.2 installed and connected to manager at
    \texttt{192.168.10.184:1514} (Agent ID 001).
  \item \textbf{Homepage migrated:} Dashboard moved to pve3 LXC 106 (Docker), fronted by
    Nginx Proxy Manager with HTTP Basic Auth (\texttt{homepage.kylemason.org}).
    Old pve1 instance (\texttt{homepage.service}) stopped and disabled.
\end{itemize}

\textbf{Deferred:} PBS backup jobs for LXC 103/104 deferred — node is temporary pending
EliteDesk G4 800 replacements. Wazuh API password reset deferred to maintenance window
(VM filesystem was repaired after corruption during reset attempts; \texttt{e2fsck -y} via
\texttt{qemu-nbd}, clean two-pass fsck confirmed).

% ── F-02 ─────────────────────────────────────────────────────────────────────
\subsection{F-02 — UPS Management Card: Telnet/FTP Exposed \quad \critical \quad \resolved}

\textbf{Host:} \texttt{192.168.10.180}

\textbf{Device (corrected):} Middle Atlantic UPS-OL2200R, IPCARD v0.970, MAPv03. This device
feeds the Dell R730s, Randy (SuperMicro), and the NetApp DS4246 JBOD. It was initially
misidentified as an APC NMC2 — the correct identification was made via TLS certificate
inspection and Telnet banner analysis.

\textbf{Description:} Telnet (port 23) and FTP (port 21) were enabled on the network
management card of the UPS feeding the primary compute stack. These protocols transmit
credentials in cleartext. A successful login to this card provides the ability to schedule
or trigger power events (graceful shutdown or immediate power cut) on attached outlets.

\textbf{Baseline:}
\begin{lstlisting}
21/tcp  open   ftp
22/tcp  closed ssh
23/tcp  open   telnet    "APC PDU/UPS devices"
80/tcp  open   http
443/tcp closed https
\end{lstlisting}

\textbf{Remediation Applied} (via authenticated menu session over Telnet, then SSH):
\begin{enumerate}[nosep]
  \item FTP disabled: \texttt{Network Settings $\to$ FTP $\to$ FTP Access $\to$ Disable}
  \item Console switched Telnet $\to$ SSH: \texttt{Network Settings $\to$ Console $\to$ Access $\to$ Enable SSH Service} (SSH host key was pre-existing and valid)
  \item Web switched HTTP $\to$ HTTPS-only: \texttt{Network Settings $\to$ Web $\to$ Access $\to$ Enable Https Service} (TLS certificate pre-existing and valid)
  \item Management card reboot via GUI to apply all pending settings (no power impact)
\end{enumerate}

\textbf{Important operational note:} Each network setting change on this card triggers a
``reflash'' of the network stack, which can silently drop other services (SSH went offline
after the web change). A full card reboot restores all services. Card reboot has zero impact
on UPS power output — the power circuitry is independent of the management card.

\textbf{Final state:}
\begin{lstlisting}
21/tcp  closed ftp       (disabled)
22/tcp  open   ssh       (restored after card reboot)
23/tcp  closed telnet    (disabled)
80/tcp  open   http      (redirect-only: 303 -> https://)
443/tcp open   https
\end{lstlisting}

SSH connection confirmed: KEX \texttt{diffie-hellman-group-exchange-sha1}, cipher
\texttt{aes128-cbc} (2010 firmware — legacy algorithms required).

% ── F-03 ─────────────────────────────────────────────────────────────────────
\subsection{F-03 — Prometheus API Unauthenticated \quad \high \quad \resolved}

\textbf{Host:} \texttt{192.168.10.183} (pve3 LXC 103)

\textbf{Description:} Prometheus (\texttt{:9090}) and Loki (\texttt{:3100}) were bound to
\texttt{0.0.0.0} in a Docker Compose stack, exposing the full infrastructure inventory
(hostnames, IP addresses, RAM, service health states) without authentication to any host on
the management VLAN. This was confirmed during the assessment:

\begin{lstlisting}
$ curl http://192.168.10.183:9090/api/v1/targets
{"job":"proxmox-randy","__address__":"192.168.10.187:9100","health":"up"}
{"job":"proxmox-pve3","__address__":"192.168.10.201:9100","health":"up"}
\end{lstlisting}

\textbf{Root cause:} Grafana was configured to reach Prometheus and Loki via the host IP
(\texttt{http://192.168.10.183:9090}) rather than Docker service names, necessitating the
ports be exposed. This was corrected as part of the remediation.

\textbf{Remediation Applied:}
\begin{enumerate}[nosep]
  \item Grafana stopped; datasource URLs updated in \texttt{grafana.db} from host IPs to
    Docker service names: \texttt{http://prometheus:9090} and \texttt{http://loki:3100}.
  \item Port bindings changed in \texttt{docker-compose.yml}:
    \texttt{'9090:9090'} $\to$ \texttt{'127.0.0.1:9090:9090'} (and \texttt{3100} likewise).
  \item Stack recreated via \texttt{docker compose up -d}. Grafana retained \texttt{0.0.0.0:3000} (has auth).
\end{enumerate}

\textbf{Verification:} Prometheus unreachable from network (000); Grafana health confirmed
(Prometheus Healthy, 3 targets up, Loki ready). Prometheus kept scraping throughout; only
Grafana had a $\sim$30s restart.

% ── F-04 ─────────────────────────────────────────────────────────────────────
\subsection{F-04 — Homepage Dashboard: No Authentication \quad \high \quad \resolved}

\textbf{Host:} \texttt{192.168.10.174} (pve1 LXC 104, moved to pve3 LXC 106)

\textbf{Description:} A Homepage (gethomepage.dev) dashboard was running as a native
Node.js process on pve1 LXC 104, bound to all interfaces on port 3000, with no
authentication. The dashboard exposed the full internal service map including URLs, ports,
and live health status of all infrastructure services.

\textbf{Remediation Applied:}
\begin{enumerate}[nosep]
  \item New dedicated LXC 106 ``homepage'' created on pve3 (Debian 12, Docker, \texttt{192.168.10.148}).
  \item Homepage config (7 YAML files, 40K) migrated from pve1 \texttt{/opt/homepage/config/} to LXC 106.
  \item Homepage deployed as official Docker image with \texttt{HOMEPAGE\_ALLOWED\_HOSTS=homepage.kylemason.org}.
  \item NPM Access List \texttt{homepage-auth} (HTTP Basic Auth, user \texttt{kyle}) created via NPM API.
  \item NPM Proxy Host \texttt{homepage.kylemason.org} created, forwarding to \texttt{192.168.10.148:3000} with the Access List attached.
  \item DOCKER-USER iptables rule applied to LXC 106: allow only NPM (\texttt{192.168.10.181}), DROP all others to port 3000. Persisted via \texttt{homepage-fw.service}.
  \item Old pve1 \texttt{homepage.service} stopped and disabled.
\end{enumerate}

\textbf{Verification:}
\begin{lstlisting}
Direct .148:3000 from network  : 000 (blocked by DOCKER-USER)
Old pve1 .174:3000             : 000 (service stopped)
homepage.kylemason.org, no auth: 401
homepage.kylemason.org, authed : 200
\end{lstlisting}

\textbf{Note:} \texttt{*.kylemason.org} resolves via public DNS (Cloudflare) to
\texttt{192.168.10.181}. A public A-record for \texttt{homepage.kylemason.org} is required
to match the existing pattern for other proxied services.

% ── F-05 ─────────────────────────────────────────────────────────────────────
\subsection{F-05 — NPM Admin Panel Exposed on Management Network \quad \medium \quad \resolved}

\textbf{Host:} \texttt{192.168.10.181:81} (pve3 LXC 101)

\textbf{Description:} The Nginx Proxy Manager admin interface (port 81) was accessible to
any host on the management VLAN. NPM controls reverse-proxy routing for all
\texttt{*.kylemason.org} services (Grafana, Vault, Wazuh, Homepage). Admin compromise would
allow route interception or service disruption across all proxied services.

\textbf{Remediation approach:} An OPNsense LAN firewall rule was initially attempted.
However, NPM (LXC 101) and all cluster nodes share the \texttt{192.168.10.0/24} L2 segment.
Inter-host traffic routes at Layer 2 via the EX3400 switch and never passes through
OPNsense. OPNsense LAN rules are ineffective for same-subnet traffic. The correct approach
is enforcement at the container level.

\textbf{Remediation Applied:} DOCKER-USER iptables rules in pve3 LXC 101, matching the
pattern established for Homepage (LXC 106):

\begin{lstlisting}
Chain DOCKER-USER:
  1. RETURN  tcp src 192.168.10.199 (Ares) dpt:81   -- allow admin workstation
  2. DROP    tcp any                dpt:81            -- block all others
\end{lstlisting}

Persisted via \texttt{/etc/systemd/system/npm-fw.service} $\to$
\texttt{/usr/local/bin/npm-fw.sh} (executes after \texttt{docker.service}).

\textbf{Verification:}
\begin{lstlisting}
Ares (.199) -> NPM :81 : 200  (allowed)
pve3 (.201) -> NPM :81 : 000  (blocked)
pve2 (.204) -> NPM :81 : 000  (blocked)
NPM :80/:443 proxy traffic    : unaffected
Grafana via NPM               : 302 (operational)
\end{lstlisting}

\textbf{Action required:} Two OPNsense LAN rules added during testing should be removed
(\texttt{Firewall $\to$ Rules $\to$ LAN}): ``Allow Ares to NPM admin'' and ``Block NPM admin
from non-Ares hosts''. These rules have no effect on L2 traffic and are misleading.

% ── F-06 ─────────────────────────────────────────────────────────────────────
\subsection{F-06 — Weak 1024-bit RSA SSH Host Key \quad \medium \quad \riskaccepted}

\textbf{Host:} \texttt{192.168.10.176}

\textbf{Device (identified):} Ubiquiti UniFi access switch. Managed by UniFi Network
controller at \texttt{192.168.10.2}. Runs Dropbear SSH 2022.83. No web UI (controller-managed).

\textbf{Description:} SSH host key is RSA $\sim$1039-bit
(\texttt{SHA256:KaROg4iOD9zcPyN+51ov15SyKpo7RsFKpFV+WXpRsqQ}), below the 2048-bit minimum.
A forged key could enable SSH MITM of the switch management session.

\textbf{Remediation Attempted:} Firmware update via the UniFi controller was completed.
UniFi switches persist the SSH host key across firmware updates — the key was unchanged.
The definitive fix (Forget + re-adopt in the controller) would reset all port configuration
and disrupt wired clients during re-provisioning.

\textbf{Risk Acceptance Rationale:}
\begin{itemize}[nosep]
  \item Controller-managed device — SSH is rarely used directly
  \item On-path attacker position required to exploit
  \item Disruption of forget + re-adopt disproportionate to risk
\end{itemize}

\textbf{Future resolution:} Forget + re-adopt in the UniFi controller, then
\texttt{ssh-keygen -R 192.168.10.176} on Ares. Schedule during a maintenance window.

% ── F-07 ─────────────────────────────────────────────────────────────────────
\subsection{F-07 — Consumer Device on Management VLAN \quad \medium \quad \pendingvlan}

\textbf{Host:} \texttt{192.168.10.112}

\textbf{Device (identified):} Amazon Fire Tablet (confirmed by owner). MAC
\texttt{XX:XX:XX:XX:XX:XX} (Amazon Technologies). mDNS hostname \texttt{linux.local}.
Running Spotify Connect (port 4070), Transmission BitTorrent (port 9091), and a SOCKS
proxy (port 1080, likely from a VPN application). Fire OS (Linux/Android kernel, TTL=64).

\textbf{Risk:} Personal consumer device with active network services on the management VLAN
intended for infrastructure only. SOCKS proxy is a potential pivot point if compromised.

\textbf{Resolution:} Move to VLAN 40 (IoT, \texttt{192.168.40.0/24}) via the UniFi
controller (reassign to IoT SSID by MAC). Dependent on VLAN segmentation activation
(pve2 trunk to EX3400 — pending infrastructure item).

% ── F-09 ─────────────────────────────────────────────────────────────────────
\subsection{F-09 --- Unidentified Devices: Amazon Echo Family \quad \low \quad \pendingvlan}

\textbf{Hosts:} \texttt{192.168.10.104--108}

\textbf{Identified:} All five hosts confirmed as Amazon Echo / Alexa family smart home
devices (Amazon Technologies OUI on all MACs). \texttt{192.168.10.108} confirmed as Amazon
Echo Show via mDNS hostname \texttt{echoshow.local}. No open TCP ports on any device;
Amazon Echo devices only initiate outbound connections to AWS services.

\textbf{Risk:} Low. Personal IoT devices on management VLAN. No inbound attack surface;
risk is network segment hygiene.

\textbf{Resolution:} Reassign all five to the IoT SSID (VLAN 40) in the UniFi controller
by MAC address once VLAN segmentation is active.

% ── F-10 ─────────────────────────────────────────────────────────────────────
\subsection{F-10 --- Unidentified Device: BBL Smart Power Strip \quad \low \quad \pendingvlan}

\textbf{Host:} \texttt{192.168.10.110}

\textbf{Identified:} BBL Technologies Co., Ltd (CN) smart power strip running an ESP32/ESP8266
WiFi SoC (Espressif MAC \texttt{XX:XX:XX:XX:XX:XX}). Confirmed via TLS certificate:
\texttt{issuer=C=CN, O=BBL Technologies Co., Ltd}; device serial
\texttt{CN=01P00C492600557}. Ports: 990 (SSL FTP), 3000 (JSON management API), 6000 (SSL).
Common default credentials rejected — device uses a custom password.

\textbf{Risk:} Low. Consumer IoT device on management VLAN. Custom credentials in place;
no default credential exposure. Network segment hygiene issue only.

\textbf{Note:} This device is \textbf{not} the APC AP7901 PDU documented in the rack
inventory (MAC \texttt{XX:XX:XX:XX:XX:XX}). The AP7901 was not visible in any network scan
and may be powered off or on a different VLAN.

\textbf{Resolution:} Identify physically (look for BBL-branded smart strip, serial
\texttt{01P00C492600557}) and reassign to VLAN 40 (IoT) via the UniFi controller.

% ─────────────────────────────────────────────────────────────────────────────
\section{Infrastructure Work Performed (Concurrent)}

The following infrastructure work was performed during the same engagement window, separate
from the penetration test remediation:

\subsection{Cluster-Wide Proxmox Upgrade}

All seven cluster nodes upgraded from various PVE 9.1.x versions to \textbf{PVE 9.2.3,
kernel 7.0.12-1-pve}. QuarkyLab retained kernel 6.14.11-9-pve (pinned) to preserve NVIDIA
550.163.01 driver compatibility (RTX 6000 24GB, the researcher's ML node).

\begin{longtable}{@{}lllll@{}}
\toprule
\textbf{Node} & \textbf{PVE Before} & \textbf{PVE After} & \textbf{Kernel} & \textbf{Notes} \\
\midrule
\endhead
pve2 & 9.1.9 & 9.2.3 & 7.0.12-1 & OPNsense autostarted \\
pve3 & 9.1.9 & 9.2.3 & 7.0.12-1 & All 4 LXCs autostarted \\
pve4 & 9.1.9 & 9.2.3 & 7.0.12-1 & \\
pve5 & 9.1.1 & 9.2.3 & 7.0.12-1 & Tailscale DNS fixed \\
Randy & 9.1.1 & 9.1.1* & 7.0.12-1 & Kernel/ZFS only \\
Jarvis & 9.1.1 & 9.2.3 & 7.0.12-1 & Root disk 6GB$\to$56GB \\
QuarkyLab & 9.1.1 & 9.2.3 & 6.14.11 (pinned) & GPU verified post-upgrade \\
\bottomrule
\end{longtable}

*Randy \texttt{pve-manager} package had no upgrade available in its repository.

\textbf{Issues encountered:}
\begin{itemize}[nosep]
  \item \textbf{Tailscale DNS bug:} Tailscale overwrites \texttt{/etc/resolv.conf} with
    MagicDNS (\texttt{100.100.100.100}) on nodes joined to Headscale, breaking apt.
    Fixed on pve4, pve5, and Jarvis: \texttt{tailscale set --accept-dns=false}.
  \item \textbf{Randy corosync split:} After first post-upgrade reboot, Randy formed a
    singleton TOTEM ring due to knet traffic disruption. Fixed via \texttt{pvecm delnode
    Randy} (from pve2), then \texttt{pkill pmxcfs; systemctl start pve-cluster} on Randy.
  \item \textbf{Jarvis root disk full:} 6GB LVM root filled during kernel package extraction
    (125MB kernel). Fixed by adding sda (186GB SAS SSD) to pve VG:
    \texttt{pvcreate/vgextend/lvextend -r}. Root now 56GB (47GB free).
  \item \textbf{pve3 LXC onboot:} All four LXCs on pve3 had \texttt{onboot=0} — would not
    restart after reboot. Set to 1 via \texttt{pct set \{101,102,103,105\} --onboot 1} before rebooting pve3.
\end{itemize}

\subsection{Headscale VPN Migration (Phase 1)}

All cluster nodes migrated from commercial Tailscale to self-hosted Headscale
(\texttt{192.168.10.186}, pve3 LXC 105, v0.29.1):

\begin{center}
\begin{tabular}{lll}
\toprule
\textbf{Node} & \textbf{Headscale IP} & \textbf{Status} \\
\midrule
Ares & 100.64.0.1 & Pre-existing \\
Randy & 100.64.0.2 & Pre-existing \\
pve5 & 100.64.0.3 & Migrated 2026-06-22 \\
pve4 & 100.64.0.4 & Migrated 2026-06-22 \\
pve3 & 100.64.0.5 & Migrated 2026-06-22 \\
Jarvis & 100.64.0.6 & Migrated 2026-06-22 \\
\bottomrule
\end{tabular}
\end{center}

\textbf{Phase 2 pending:} QuarkyLab and the researcher's Mac must migrate together — the
researcher uses QuarkyLab via Tailscale for ML work (device/user identifiers kept in the
private ops vault). Coordinate with the researcher before migrating.

\subsection{PBS Backup Schedules Configured}

Proxmox Backup Server (Randy, \texttt{https://192.168.10.187:8007}) now has active backup
jobs for all guest workloads:

\begin{center}
\begin{tabular}{lll}
\toprule
\textbf{Job} & \textbf{Schedule} & \textbf{Retention} \\
\midrule
LXCs (101, 102, 103, 105) & Daily 02:00 & 7 daily, 4 weekly \\
VMs (100, 104) & Daily 03:00 & 7 daily, 4 weekly \\
\bottomrule
\end{tabular}
\end{center}

Test backup of LXC 105 (Headscale) completed successfully in 5 seconds. PBS datastore
confirmed at \texttt{/datastore/ct/105}.

% ─────────────────────────────────────────────────────────────────────────────
\section{Remediation Checklist}

\begin{longtable}{@{}p{0.8cm}p{7.5cm}p{2cm}p{2cm}@{}}
\toprule
\textbf{ID} & \textbf{Action} & \textbf{Owner} & \textbf{Status} \\
\midrule
\endhead
F-01 & Harden pve1 SSH / enroll Wazuh agent & System & \resolved \\
F-01 & PBS backups for pve1 LXCs 103+104 & System & Deferred \\
F-01 & Decommission pve1 when EliteDesk G4 800s arrive & Hardware & Pending \\
F-02 & Disable Telnet + FTP on UPS NMC & System & \resolved \\
F-02 & Enable SSH + HTTPS on UPS NMC & System & \resolved \\
F-03 & Bind Prometheus + Loki to localhost & System & \resolved \\
F-04 & Migrate Homepage to pve3 + NPM auth & System & \resolved \\
F-05 & Restrict NPM :81 to Ares via DOCKER-USER & System & \resolved \\
F-05 & Remove redundant OPNsense LAN rules & Admin & Pending \\
F-06 & Firmware update (key unchanged) & System & Done \\
F-06 & Forget + re-adopt UniFi switch & Admin & Deferred \\
F-07 & Move Fire Tablet to VLAN 40 & Admin & Pending VLAN \\
F-09 & Move Echo devices to VLAN 40 & Admin & Pending VLAN \\
F-10 & Identify BBL strip physically + move to VLAN 40 & Admin & Pending VLAN \\
— & Activate VLAN segmentation (pve2 trunk) & Infrastructure & Pending \\
— & Add public DNS A: \texttt{homepage.kylemason.org} $\to$ \texttt{192.168.10.181} & DNS & Pending \\
— & Locate APC AP7901 PDU (MAC \texttt{XX:XX:XX:XX:XX:XX}) & Physical & Pending \\
— & Wazuh API password reset (maintenance window) & System & Deferred \\
— & Headscale Phase 2: QuarkyLab + the researcher's Mac & Admin & Pending \\
\bottomrule
\end{longtable}

% ─────────────────────────────────────────────────────────────────────────────
\section{Technical Notes}

\subsection{OPNsense Firewall and L2 Traffic}

A key architectural finding during F-05 remediation: OPNsense firewall rules on the LAN
interface only intercept traffic that \textit{routes through} OPNsense (cross-subnet or
WAN-bound traffic). All hosts on \texttt{192.168.10.0/24} communicate at Layer 2 via the
Juniper EX3400 switch. Traceroutes between any two management-VLAN hosts show a single hop
directly to the destination. \textbf{OPNsense LAN rules cannot restrict same-subnet traffic.}

The correct enforcement mechanism for same-subnet access control is at the container or
host level (iptables DOCKER-USER chain, as implemented for F-04 and F-05).

\subsection{Tailscale DNS Interaction with apt}

Nodes enrolled in Headscale with Tailscale installed receive a rewritten
\texttt{/etc/resolv.conf} pointing to MagicDNS (\texttt{100.100.100.100}). Since MagicDNS
is not configured on the self-hosted Headscale instance, DNS resolution fails silently —
causing \texttt{apt-get} to hang indefinitely when attempting to fetch package lists.

\textbf{Permanent fix} for any Headscale-enrolled node before apt operations:
\begin{lstlisting}
tailscale set --accept-dns=false
echo 'nameserver 192.168.10.177' > /etc/resolv.conf
echo 'nameserver 192.168.10.1' >> /etc/resolv.conf
\end{lstlisting}

\subsection{Middle Atlantic UPS NMC: Single-Session Limit}

The UPS management card enforces a single concurrent management session. If a Telnet or SSH
session terminates uncleanly (without using menu option ``Logout''), the slot remains locked
for approximately 5 minutes. Subsequent connection attempts receive an immediate reset.
Always log out cleanly via the menu.

% ─────────────────────────────────────────────────────────────────────────────
\section{Appendix A: Network Asset Inventory (Updated)}

\begin{longtable}{@{}llll@{}}
\toprule
\textbf{IP} & \textbf{Identity} & \textbf{MAC} & \textbf{VLAN} \\
\midrule
\endhead
192.168.10.1 & OPNsense (VM 100, pve2) & XX:XX:XX:XX:XX:XX & 1 \\
192.168.10.2 & UniFi Network Controller & XX:XX:XX:XX:XX:XX & 1 \\
192.168.10.20 & QuarkyLab iDRAC & XX:XX:XX:XX:XX:XX & 1 \\
192.168.10.21 & Jarvis iDRAC & XX:XX:XX:XX:XX:XX & 1 \\
192.168.10.22 & Randy IPMI & XX:XX:XX:XX:XX:XX & 1 \\
192.168.10.31 & Jarvis (Dell R730) & XX:XX:XX:XX:XX:XX & 1 \\
192.168.10.50 & Juniper EX3400-48P & XX:XX:XX:XX:XX:XX & 1 \\
192.168.10.100 & Ares (admin workstation) & XX:XX:XX:XX:XX:XX & 1 \\
192.168.10.104--108 & Amazon Echo devices (×5) & 08:C2:24:* / 98:CC:F3:* & Should be VLAN 40 \\
192.168.10.110 & BBL smart power strip & XX:XX:XX:XX:XX:XX & Should be VLAN 40 \\
192.168.10.111 & Sonos device & — & 1 \\
192.168.10.112 & Amazon Fire Tablet & XX:XX:XX:XX:XX:XX & Should be VLAN 40 \\
192.168.10.122 & SpotifyConnect device & XX:XX:XX:XX:XX:XX & 1 \\
192.168.10.148 & Homepage (pve3 LXC 106) & BC:24:11:* & 1 \\
192.168.10.172 & Brother MFC-J4335DW printer & XX:XX:XX:XX:XX:XX & 1 \\
192.168.10.174 & pve1 LXC 104 (Homepage — stopped) & XX:XX:XX:XX:XX:XX & 1 \\
192.168.10.176 & UniFi access switch & XX:XX:XX:XX:XX:XX & 1 \\
192.168.10.177 & Pi-hole (pve1 LXC 103) & XX:XX:XX:XX:XX:XX & 1 \\
192.168.10.179 & QuarkyLab (Dell R730) & — & 1 \\
192.168.10.180 & Middle Atlantic UPS-OL2200R NMC & XX:XX:XX:XX:XX:XX & 1 \\
192.168.10.181 & Nginx Proxy Manager (pve3 LXC 101) & XX:XX:XX:XX:XX:XX & 1 \\
192.168.10.182 & Vaultwarden (pve3 LXC 102) & XX:XX:XX:XX:XX:XX & 1 \\
192.168.10.183 & Prometheus/Grafana/Loki (pve3 LXC 103) & XX:XX:XX:XX:XX:XX & 1 \\
192.168.10.184 & Wazuh SIEM (QuarkyLab VM 104) & XX:XX:XX:XX:XX:XX & 1 \\
192.168.10.186 & Headscale VPN (pve3 LXC 105) & XX:XX:XX:XX:XX:XX & 1 \\
192.168.10.187 & Randy (SuperMicro, PBS) & XX:XX:XX:XX:XX:XX & 1 \\
192.168.10.191 & Secondary laptop & — & 1 \\
192.168.10.193 & pve1 Mac Mini (standalone) & XX:XX:XX:XX:XX:XX & 1 \\
192.168.10.198 & SpotifyConnect device & — & 1 \\
192.168.10.199 & Ares (admin workstation, wired) & — & 1 \\
192.168.10.201 & pve3 (HP EliteDesk G4) & XX:XX:XX:XX:XX:XX & 1 \\
192.168.10.202 & pve4 (HP EliteDesk G3 Mini) & XX:XX:XX:XX:XX:XX & 1 \\
192.168.10.203 & pve5 (HP EliteDesk G3 Mini) & XX:XX:XX:XX:XX:XX & 1 \\
192.168.10.204 & pve2 (HP EliteDesk G4 SFF) & XX:XX:XX:XX:XX:XX & 1 \\
\bottomrule
\end{longtable}

% ─────────────────────────────────────────────────────────────────────────────
\section{Appendix B: File Locations}

\begin{tabular}{ll}
\toprule
\textbf{Item} & \textbf{Path} \\
\midrule
This report (source) & \texttt{Home-Lab/docs/netframe-pentest-remediation-2026-06-24.tex} \\
Remediation log & \texttt{Home-Lab/docs/pentest-remediation-log.md} \\
Original pentest report & \texttt{Home-Lab/docs/pentest-2026-06-23.md} \\
NPM firewall script & \texttt{pve3 LXC 101:/usr/local/bin/npm-fw.sh} \\
Homepage firewall script & \texttt{pve3 LXC 106:/usr/local/bin/homepage-fw.sh} \\
Homepage compose & \texttt{pve3 LXC 106:/opt/homepage/docker-compose.yml} \\
Prometheus compose & \texttt{pve3 LXC 103:/opt/grafana/docker-compose.yml} \\
step-ca password & stored in Vaultwarden (path not published) \\
\bottomrule
\end{tabular}

\vfill
\begin{center}
\rule{0.6\linewidth}{0.4pt}\\[0.5em]
{\small\color{gray}
NetFRAME Homelab $\cdot$ Kyle Mason $\cdot$ Greater Cleveland, OH $\cdot$ 2026-06-24\\
CONFIDENTIAL — For personal use only
}
\end{center}

\end{document}
